% ========== 1.5 RELATED WORK ==========

\section{1.6 Related Work}
\addcontentsline{toc}{section}{1.6 Related Work}
\label{sec:related-work}

\lettrine{T}raditional integrated development environments like Visual Studio, Eclipse, and IntelliJ IDEA established foundational concepts for modern development tools. These systems demonstrated the value of integrating multiple development tools into unified environments. They introduced plugin architectures that enabled community-driven extensibility. However, their monolithic designs and human-centric architectures create barriers when attempting to integrate sophisticated AI capabilities.

Modern code editors represent an evolution toward lighter-weight tools. Visual Studio Code achieved market dominance through balanced performance and extensive extension marketplace. Sublime Text demonstrated that native performance could coexist with extensibility. These editors improved on traditional IDEs but maintained architectural patterns designed for human-centric workflows rather than AI collaboration.

Recent AI-assisted development tools mark a significant shift in capabilities. GitHub Copilot introduced mainstream developers to AI-powered code completion using large language models. ChatGPT enabled conversational interfaces for debugging and code explanation. Tools like Cursor and Windsurf attempt deeper AI integration through custom implementations. However, these tools still operate primarily as assistants rather than collaborative partners, and they suffer from the performance and integration limitations of retrofit approaches.

Symphony differs from existing work by designing for AI collaboration from the ground up rather than adding AI features to traditional architectures. Where current tools treat AI as supplementary features, Symphony positions AI as a foundational architectural principle. This fundamental difference enables capabilities that would be impossible to achieve through incremental improvements to existing systems.
